\documentclass[11pt,a4paper]{article}
\usepackage[margin=25mm]{geometry}
\usepackage{iftex}
\ifPDFTeX
  \usepackage[T2A]{fontenc}
  \usepackage[utf8]{inputenc}
\else
  \usepackage{fontspec}
  \setmainfont{cmunrm.otf}[
    BoldFont=cmunbx.otf,
    ItalicFont=cmunti.otf,
    BoldItalicFont=cmunbi.otf,
    Ligatures=TeX]
\fi
\usepackage[english]{babel}
\usepackage{amsmath,amssymb,amsthm}
\usepackage{microtype}
\usepackage{enumitem}
\usepackage[unicode]{hyperref}
\hypersetup{
  pdftitle={Extensions of Quandles with Stabilizing Orbit Series},
  pdfauthor={Danil Pokulevskii},
  pdflang={en-US},
  colorlinks=true,
  linkcolor=blue,
  citecolor=blue,
  urlcolor=blue}
\setlength{\emergencystretch}{2em}
\setlist[enumerate]{itemsep=2pt,topsep=5pt}
\frenchspacing

\newtheorem{theorem}{Theorem}[section]
\newtheorem{lemma}[theorem]{Lemma}
\newtheorem{proposition}[theorem]{Proposition}
\newtheorem{corollary}[theorem]{Corollary}
\theoremstyle{definition}
\newtheorem{definition}[theorem]{Definition}
\theoremstyle{remark}
\newtheorem{remark}[theorem]{Remark}
\renewcommand{\proofname}{Proof}
\newcommand{\OS}{\mathcal{OS}}
\newcommand{\tOS}{t\mathcal{OS}}
\newcommand{\Inn}{\operatorname{Inn}}
\newcommand{\Orb}{\operatorname{Orb}}
\newcommand{\Aut}{\operatorname{Aut}}
\newcommand{\Z}{\mathbb Z}
\newcommand{\conj}{\triangleright}

\title{Extensions of Quandles\\with Stabilizing Orbit Series}
\author{Danil Pokulevskii\\[3pt]
  \small Novosibirsk State Technical University\\
  \small Novosibirsk, Russia\\
  \small\href{mailto:555tery@gmail.com}{\texttt{555tery@gmail.com}}}
\date{October 8, 2026\\[2pt]\small Working preprint}

\begin{document}
\maketitle

\begin{abstract}
Let the class \(\OS\) consist of quandles in which every orbit
series stabilizes after finitely many steps. We prove that
\(\OS\) is closed under extensions: if the quotient and every full
fibre of an epimorphism belong to \(\OS\), then the original quandle
does as well. No common bound on the depths of the fibres is required.
The key ingredient is a group identity for relative commutator
series of a single automorphism. If \(N\trianglelefteq L\),
\(\phi(N)=N\), and \([L/N,\bar\phi]=L/N\), then the series
\(R_0=N\), \(R_{i+1}=[R_i,\phi]\) and
\(D_0=L\), \(D_{i+1}=[D_i,\phi]\) satisfy
\(D_i=R_iD_{i+1}\). The passage to quandles proceeds through
conjugacy classes and coverings and does not require injectivity of the map
that assigns to each element its left translation. We also show that the class
\(\OS_\omega=\bigcup_{n\ge1}\OS_n\) is not closed under extensions.
Thus we answer both parts of Question~5.7 of
Bonatto--Crans--Nasybullov--Whitney.
\end{abstract}

\noindent\textbf{Keywords:} quandle, orbit series, extension,
subnormal subgroup, group automorphism.

\section{Introduction and main results}

Orbit series of quandles, the classes \(\OS\), \(\OS_n\),
\(\tOS\), \(\tOS_n\), and their unions over finite indices
were introduced in \cite[Section~4.1]{BCNW}.
Question~5.7 of that paper asks whether
\(\OS\) and \(\OS_\omega\) are closed under extensions.
The numbering of the question and results from \cite{BCNW} used below
follows arXiv version v3, dated August 18, 2020.

It is essential to distinguish two conditions. In the class \(\OS\),
each individual orbit series must stabilize, but the number of
steps may depend on the series. In the class \(\OS_\omega\), a single finite
bound must work for all series of a given quandle.
The definition of an extension does not require a common bound across
different full fibres.

\begin{theorem}\label{thm:main}
Let \(\pi:Q\twoheadrightarrow B\) be an epimorphism of quandles.
If \(B\in\OS\) and \(\pi^{-1}(b)\in\OS\) for every \(b\in B\),
then \(Q\in\OS\). Consequently, the class \(\OS\) is closed under
extensions.
\end{theorem}

\begin{proposition}\label{prop:negative}
The class \(\OS_\omega\) is not closed under extensions.
A counterexample can be chosen with a trivial base and finite
full fibres.
\end{proposition}

Theorem~\ref{thm:main} is proved in Section~\ref{sec:quandles}.
Its group-theoretic core, Theorem~\ref{thm:group}, states that
stabilization of a relative commutator series on a normal
subgroup lifts to the whole group without increasing the number of steps,
provided the quotient is perfect relative to the automorphism
under consideration. The proof uses the elementary fact that
a subnormal subgroup of finite defect that normally generates a
group must equal that group.
Proposition~\ref{prop:negative} follows from the disjoint union of
dihedral quandles in Example~4.2 of \cite{BCNW},
after considering the projection onto the summand index.

The definitions of orbit-series classes and extensions follow~\cite{BCNW};
the term ``covering'' is used in the sense of~\cite{Eisermann},
translated to the left convention. All auxiliary group and quandle
lemmas, including the required property of coverings,
are proved below. For the negative example, we also give
a direct computation of orbit series.

\section{Definitions and auxiliary lemmas}\label{sec:prelim}

\begin{definition}
\emph{A quandle} is a nonempty set \(Q\) with an operation \(\conj\)
satisfying
\[
 x\conj x=x,\qquad
 L_x:y\longmapsto x\conj y\text{ is bijective},\qquad
 x\conj(y\conj z)=(x\conj y)\conj(x\conj z).
\]
A subquandle is closed under the operation and inverse left translations.
The group \(\Inn(Q)=\langle L_x:x\in Q\rangle\) is called the
inner automorphism group of the quandle. We write \(\Orb(x,Q)\) for
the orbit of \(x\). A quandle is connected if this group acts transitively;
it is trivial if \(x\conj y=y\) for all \(x,y\).
\end{definition}

A quandle homomorphism preserves the operation; an epimorphism
here and below means a surjective homomorphism.

By choosing \(x_i\in Q_i\) successively, we construct an \emph{orbit series}
\[
 Q_0=Q,\qquad Q_{i+1}=\Orb(x_i,Q_i).
\]
The class \(\OS\) consists of quandles such that every orbit series
has some \(m\ge0\) with \(Q_m=Q_{m+1}\).
The class \(\OS_n\), for \(n\ge1\), is defined by requiring that such an
\(m\le n\) can be chosen for every series. Set
\(\OS_\omega=\bigcup_{n\ge1}\OS_n\).
The stable term of a series is connected, since
\(Q_m=\Orb(x_m,Q_m)\).
The trivializing classes \(\tOS\) and \(\tOS_n\) are defined by
requiring that a singleton quandle be reached, respectively,
in finitely many steps or no later than step \(n\).
Set \(\tOS_\omega=\bigcup_{n\ge1}\tOS_n\).

By an \emph{extension} we mean an epimorphism
\(\pi:Q\twoheadrightarrow B\); its full fibres are
\(F_b=\pi^{-1}(b)\). Equivalently, one considers the quotient
\(Q/\alpha\) and all full congruence classes of \(\alpha=\ker\pi\),
as in \cite[Section~2.2]{BCNW}.

\begin{lemma}\label{lem:images}
An epimorphism \(p:E\twoheadrightarrow C\) induces an epimorphism
\(\Inn(E)\twoheadrightarrow\Inn(C)\) sending \(L_x\)
to \(L_{p(x)}\), and
\[
 p(\Orb(x,E))=\Orb(p(x),C).
\]
In particular, images of orbit series are orbit series,
and the class \(\OS\) is closed under epimorphic images.
\end{lemma}
\begin{proof}
The identity \(pL_x=L_{p(x)}p\) also holds for inverse translations.
If a word in translations acts trivially on \(E\), then the corresponding
word acts trivially on \(C\), since \(p\) is surjective.
This makes the group map well-defined, and its surjectivity
follows by lifting generators. The equality of orbits follows.
For an arbitrary subquandle, the same argument applies after restricting
the map to that subquandle and taking its image as the codomain.

Any orbit series in \(C\) can be lifted successively to a series
in \(E\): each next anchor can be lifted to the current term, which
maps surjectively onto the corresponding term in \(C\).
Stabilization of the lifted series implies stabilization of the original one.
\end{proof}

\begin{lemma}\label{lem:invariant}
If \(F\in\OS\) and a nonempty subquandle \(S\subseteq F\)
is invariant under \(\Inn(F)\), then \(S\in\OS\).
\end{lemma}
\begin{proof}
Take any series \(S=S_0\supseteq S_1\supseteq\cdots\)
with anchors \(x_j\in S_j\). Using the same anchors, define
\(F_0=F\), \(F_{j+1}=\Orb(x_j,F_j)\).
We obtain by induction
\[
 F_j\supseteq S_j\supseteq F_{j+1},
 \qquad S_j\text{ is invariant under }\Inn(F_j).
\]
Indeed, invariance of \(S_j\) implies
\(F_{j+1}\subseteq S_j\), while \(S_j\subseteq F_j\) implies
\(S_{j+1}\subseteq F_{j+1}\).
Since \(S_{j+1}\) is an orbit in \(S_j\), and
\(F_{j+1}\subseteq S_j\), translations by elements of \(F_{j+1}\)
and their inverses preserve \(S_{j+1}\). This gives the next step.

The series \(F_j\) stabilizes:
\(F_m=F_{m+1}=\Orb(x_m,F_m)\).
The inclusions give \(S_m=F_m\); connectedness of \(F_m\) implies
\(S_{m+1}=S_m\).
\end{proof}

\begin{remark}
In Lemma~\ref{lem:invariant}, the invariance assumption is essential.
We do not use inheritance of \(\OS\) by arbitrary subquandles;
in general, this fails by \cite[Example~5.3]{BCNW}.
\end{remark}

\begin{lemma}\label{lem:slices}
Suppose all full fibres of an epimorphism
\(\pi:Q\twoheadrightarrow B\) belong to \(\OS\),
and let \(Q_i\) be an orbit series in \(Q\).
Then every nonempty slice \(Q_i\cap F_b\) belongs to \(\OS\).
\end{lemma}
\begin{proof}
If \(S_i=Q_i\cap F_b\ne\varnothing\), then all
\(S_j=Q_j\cap F_b\), \(0\le j\le i\), are nonempty.
For \(j<i\) and \(y\in S_j\), the translations \(L_y^{\pm1}\)
preserve both the orbit \(Q_{j+1}\) in \(Q_j\) and the fibre \(F_b\):
on the base, \(L_b^{\pm1}\) fix \(b\).
Thus \(S_{j+1}\) is invariant under \(\Inn(S_j)\).
Starting with \(S_0=F_b\), we successively apply
Lemma~\ref{lem:invariant} up to index \(i\).
Nonemptiness of later slices is not required.
\end{proof}

\begin{lemma}\label{lem:trivialbase}
An epimorphism onto a trivial quandle with full fibres in \(\OS\)
has its domain in \(\OS\).
\end{lemma}
\begin{proof}
The first orbit of any series lies in a single full fibre.
As an orbit under the inner group of the domain, it is invariant under
the inner group of that fibre. Apply
Lemma~\ref{lem:invariant}; the remaining series stabilizes.
\end{proof}

\begin{definition}
An epimorphism \(p:E\twoheadrightarrow C\) is called a
\emph{covering} if
\[
p(x)=p(y)\quad\Longrightarrow\quad L_x^E=L_y^E.
\]
This is \cite[Definition~2.32, arXiv v3]{Eisermann},
translated from the right convention to the left convention.
\end{definition}

\begin{lemma}\label{lem:covering}
If \(p:E\twoheadrightarrow C\) is a covering and \(C\in\OS\),
then \(E\in\OS\).
\end{lemma}
\begin{proof}
For a series \(E_i\), the images \(C_i=p(E_i)\) stabilize
at a connected subquandle \(C_*\). On the corresponding tail, all
\(E_i\twoheadrightarrow C_*\).
Consider the groups of global permutations of the set \(E\)
\[
 \Gamma_i=\langle L_x^E:x\in E_i\rangle.
\]
Their generating sets are the same on the tail: a translation
depends only on \(p(x)\), and the image sets \(p(E_i)=C_*\)
coincide. Denote the common group by \(\Gamma\).
Each \(E_i\) is preserved by this group, and its restrictions
generate \(\Inn(E_i)\).
Therefore, \(E_{i+1}\) is a \(\Gamma\)-orbit;
the next anchor lies in this same orbit, and
\(E_{i+2}=E_{i+1}\).
\end{proof}

\section{Lifting a relative commutator series}\label{sec:groups}

For a group, write \(x^y=y^{-1}xy\).
The ordinary commutator is \([x,y]=x^{-1}y^{-1}xy\),
\(H'=[H,H]\) is the commutator subgroup,
and \(C_H(h)=\{q\in H:qh=hq\}\) is the centralizer.
The normal closure of a subgroup is the smallest normal subgroup
containing it.
If \(\phi\in\Aut(L)\) and \(X\le L\) satisfies
\(\phi(X)=X\), set
\[
 [X,\phi]=\langle x^{-1}\phi(x):x\in X\rangle.
\]
The notation denotes the subgroup generated by these elements, not merely the set of values.

\begin{lemma}\label{lem:subnormal}
Let \(A\) be subnormal in \(B\) with finite defect.
If \(\langle A^B\rangle=B\), then \(A=B\).
\end{lemma}
\begin{proof}
Choose a finite chain
\(A=A_0\trianglelefteq A_1\trianglelefteq\cdots
\trianglelefteq A_m=B\).
If \(m=0\), there is nothing to prove.
Otherwise, the normal subgroup \(A_{m-1}\) of \(B\) contains \(A\),
and hence contains its normal closure, which is \(B\).
Thus \(A_{m-1}=B\). Removing the last terms successively gives \(A=B\).
\end{proof}

\begin{theorem}\label{thm:group}
Let \(N\trianglelefteq L\), \(\phi\in\Aut(L)\),
\(\phi(N)=N\), and suppose the induced automorphism \(\bar\phi\)
satisfies \([L/N,\bar\phi]=L/N\).
Define
\[
 R_0=N,\quad R_{i+1}=[R_i,\phi],\qquad
 D_0=L,\quad D_{i+1}=[D_i,\phi].
\]
Then, for every \(i\ge0\),
\begin{equation}\label{eq:strong}
 D_i=R_iD_{i+1}.
\end{equation}
In particular,
\[
 R_r=R_{r+1}\quad\Longrightarrow\quad D_r=D_{r+1}
 \qquad(r\ge0).
\]
\end{theorem}
\begin{proof}
Write \(\delta(x)=x^{-1}\phi(x)\). We have
\begin{equation}\label{eq:delta}
 \delta(xy)=\delta(x)^y\delta(y),\qquad
 \delta(x)^y=\delta(xy)\delta(y)^{-1}.
\end{equation}
It follows that \([X,\phi]\trianglelefteq X\) for every \(\phi\)-invariant \(X\).
The identity \(\phi(\delta(x))=\delta(\phi(x))\) shows
\(\phi\)-invariance of this subgroup.
The series \(R_i,D_i\) are descending; monotonicity of the operation
\(X\mapsto[X,\phi]\) gives \(R_i\le D_i\).
Moreover, \(R_i\) is subnormal in \(D_i\):
intersect the finite chain
\[
 R_i\trianglelefteq R_{i-1}\trianglelefteq\cdots
 \trianglelefteq N\trianglelefteq L
\]
with \(D_i\). Normality of consecutive terms is preserved;
the endpoints of the resulting chain are \(R_i\) and \(D_i\).

We prove \eqref{eq:strong} by induction.
For \(i=0\), the image of \(D_1\) in \(L/N\) is the whole quotient,
so \(L=ND_1\).
Let \(i\ge1\) and suppose \(D_{i-1}=R_{i-1}D_i\).
Set \(M_i=\langle R_i^{D_i}\rangle\).
Every \(x\in D_{i-1}\) has the form \(rd\),
where \(r\in R_{i-1}\) and \(d\in D_i\).
By \eqref{eq:delta},
\[
 \delta(x)=\delta(r)^d\delta(d)\in M_iD_{i+1}.
\]
Since \(D_{i+1}\trianglelefteq D_i\), the right-hand side is a subgroup.
All generators of \(D_i=[D_{i-1},\phi]\) lie in it,
and the reverse inclusion is immediate. Thus
\[
 D_i=M_iD_{i+1}.
\]
The image of \(R_i\) in \(D_i/D_{i+1}\) is subnormal of finite defect
and normally generates this quotient.
By Lemma~\ref{lem:subnormal}, it is the entire quotient,
which gives \eqref{eq:strong} and completes the induction.

If \(R_r=R_{r+1}\), then
\(R_r=[R_r,\phi]\le[D_r,\phi]=D_{r+1}\).
Now \eqref{eq:strong} gives \(D_r=D_{r+1}\).
\end{proof}

\begin{remark}
Theorem~\ref{thm:group} assumes no finiteness,
solvability, splitting, or finite order of the automorphism.
The terms \(R_i\) need not be normal in all of \(L\);
only their subnormality in \(D_i\) is used.
\end{remark}

\section{Extensions of conjugacy classes}\label{sec:conjugacy}

We regard the class \(c^K=\{g^{-1}cg:g\in K\}\) as a quandle
with operation \(a\conj b=aba^{-1}\).
This convention corresponds to \(\mathrm{Conj}_{-1}\) in \cite{BCNW}.

\begin{lemma}\label{lem:towers}
For \(c\in K\), set \(K_0=K\),
\(K_{j+1}=\langle c^{K_j}\rangle\).
The quandle \(c^K\) belongs to \(\OS\) if and only if
the subgroup series \(K_j\) stabilizes after finitely many steps.
\end{lemma}
\begin{proof}
Set \(Y_j=c^{K_j}\).
Then \(K_{j+1}\trianglelefteq K_j\) and
\(Y_{j+1}=\Orb(c,Y_j)\), so \(Y_j\) is a principal
orbit series with constant anchor \(c\).
Equality \(K_j=K_{j+1}\) gives \(Y_j=Y_{j+1}\), while equality
\(Y_j=Y_{j+1}\) gives \(K_{j+1}=K_{j+2}\).

We show that stabilization of the principal series is sufficient.
Any finite initial segment of an arbitrary orbit
series can be conjugated into the form \(Y_j\).
If the first terms have already been put in this form, the next anchor has the form
\(c^g\), where \(g\in K_j\).
Conjugation by \(g^{-1}\) sends it to \(c\)
and preserves all previous \(Y_k\), for \(k\le j\),
since \(K_j\le K_k\).
It sends the next orbit to \(Y_{j+1}\).
Thus, if \(Y_m=Y_{m+1}\), normalization through
the anchor of index \(m\) gives equality of the corresponding
terms in any series. The converse follows from
stabilization of the principal series.
\end{proof}

\begin{lemma}\label{lem:kernelfiber}
Let \(\sigma:G\twoheadrightarrow H\), \(h=\sigma(c)\),
\(N=\ker\sigma\), \(P=\sigma^{-1}(C_H(h))\).
If the full fibre \(c^P\) of the epimorphism \(c^G\to h^H\)
belongs to \(\OS\), then \(c^N\in\OS\).
\end{lemma}
\begin{proof}
Set \(V=N\langle c\rangle=\sigma^{-1}(\langle h\rangle)\).
Since \(\langle h\rangle\) is central in \(C_H(h)\),
we have \(V\trianglelefteq P\).
Every element of \(c^P\) maps to \(h\), so
\(\langle c^P\rangle\le V\). Moreover, \(c^V=c^N\):
every element of \(V=\langle c\rangle N\) can be written as \(c^k n\),
and \(c^{c^k n}=c^n\).
Normality of \(N\) ensures both orders of decomposition of \(V\).

For \(A_0=P\), \(B_0=V\), and
\[
 A_{j+1}=\langle c^{A_j}\rangle,\qquad
 B_{j+1}=\langle c^{B_j}\rangle
\]
monotonicity gives, by induction,
\[
 A_{j+1}\le B_j\le A_j\qquad(j\ge0).
\]
The base case is \(A_1\le V\le P\).
Applying the monotone operator to the two inclusions
gives the next two inclusions.
If \(A_m=A_{m+1}=T\), then \(A_{m+2}=T\) as well;
squeezing gives \(B_m=B_{m+1}=T\).
We now apply Lemma~\ref{lem:towers}.
\end{proof}

\begin{proposition}\label{prop:groupquandle}
Let \(\sigma:G\twoheadrightarrow H\), \(h=\sigma(c)\), and
\(H=\langle h^H\rangle\).
If the full fibre \(c^P\), where \(P=\sigma^{-1}(C_H(h))\),
belongs to \(\OS\), then \(c^G\in\OS\).
\end{proposition}
\begin{proof}
Let \(H'\) be the commutator subgroup of \(H\), \(N=\ker\sigma\),
\(L=\sigma^{-1}(H')\).
Normal generation of \(H\) by \(h\) implies
\[
 H=H'\langle h\rangle,\qquad G=L\langle c\rangle.
\]
The restriction \(\sigma|_L\) maps onto \(H'\)
and \(L/N\cong H'\).
Set \(\phi(l)=c^{-1}lc\), \(\psi(q)=h^{-1}qh\).
Since \(N=\ker\sigma\), we have \(N\trianglelefteq L\)
and \(\phi(N)=N\).
The induced automorphism \(\bar\phi\) corresponds to \(\psi\).

\(H'\) is characteristic in \(H\), hence invariant under \(\psi\).
The subgroup \(J=[H',\psi]\) is normal in \(H'\)
and invariant under \(h\), hence normal in \(H\).
In \(H/J\), the image of \(h\) is central and normally generates the whole quotient.
The quotient is cyclic; hence \(H'\le J\), that is,
\begin{equation}\label{eq:headperfect}
 [H',\psi]=H'.
\end{equation}

Introduce auxiliary groups with infinite cyclic
generators \(u,v\):
\[
 \widetilde G=L\rtimes_\phi\langle u\rangle,\qquad
 \widetilde H=H'\rtimes_\psi\langle v\rangle,
 \qquad u^{-1}lu=\phi(l),\quad v^{-1}qv=\psi(q).
\]
The following epimorphisms are well-defined:
\[
 \widetilde\sigma(lu^k)=\sigma(l)v^k,\qquad
\rho_G(lu^k)=lc^k,\qquad \rho_H(qv^k)=qh^k.
\]
Indeed, \(\sigma(\phi(l))=\psi(\sigma(l))\)
and \(c^{-1}lc=\phi(l)\), \(h^{-1}qh=\psi(q)\),
so the defining relations are preserved.
Surjectivity follows from \(\sigma(L)=H'\),
\(G=L\langle c\rangle\) and \(H=H'\langle h\rangle\).
The first map has kernel \(N\), and the other maps are compatible
with \(\sigma\) and \(\widetilde\sigma\).
The restrictions of \(\rho_G,\rho_H\) induce quandle isomorphisms
\begin{equation}\label{eq:liftclasses}
 u^{\widetilde G}\cong c^G,\qquad
 v^{\widetilde H}\cong h^H.
\end{equation}
Indeed, surjectivity of the group maps allows any conjugating element
to be lifted. Every element of \(u^{\widetilde G}\)
has \(\Z\)-coordinate \(1\),
and on \(Lu\), the map \(\rho_G\) is injective:
\(lc=l'c\) implies \(l=l'\).
The argument for the base is the same.
The isomorphisms in \eqref{eq:liftclasses} are compatible with the epimorphisms
of quandles and preserve all their full fibres;
we also have \(u^N\cong c^N\).

For a \(\phi\)-invariant \(X\le L\), the normal closure of \(u\)
in \(X\rtimes_\phi\langle u\rangle\) is
\begin{equation}\label{eq:closureu}
 [X,\phi]\rtimes\langle u\rangle.
\end{equation}
It contains \(u\) and all
\(x^{-1}\phi(x)=(u^x)^{-1}u\).
Conversely, the subgroup on the right is normal:
\([X,\phi]\trianglelefteq X\), it is invariant under \(u\),
and on \(X/[X,\phi]\), the action of \(u\) is trivial.
Consequently, it contains the normal closure of \(u\).

By Lemma~\ref{lem:kernelfiber}, \(c^N\in\OS\), hence
\(u^N\in\OS\). Also,
\(u^{N\rtimes\langle u\rangle}=u^N\).
By \eqref{eq:closureu}, the series of normal closures of \(u\)
in this group has terms \(R_i\rtimes\langle u\rangle\),
where \(R_0=N\), \(R_{i+1}=[R_i,\phi]\).
Lemma~\ref{lem:towers} gives stabilization of \(R_i\).
Condition \eqref{eq:headperfect} allows us to apply
Theorem~\ref{thm:group}; the global series
\(D_0=L\), \(D_{i+1}=[D_i,\phi]\) stabilizes.
By \eqref{eq:closureu}, the series of normal
closures of \(u\) in \(\widetilde G\) stabilizes as well.
Thus \(u^{\widetilde G}\in\OS\), and
\eqref{eq:liftclasses} gives \(c^G\in\OS\).
\end{proof}

\begin{remark}
The auxiliary semidirect product is not an assumption
that the original \(\sigma\) splits.
The element \(c\) may have finite order;
the generator \(u\) is chosen to have infinite order only in the
auxiliary group.
\end{remark}

\section{Proof of extension closure for arbitrary quandles}
\label{sec:quandles}

\begin{lemma}\label{lem:connectedbase}
Let \(\pi:E\twoheadrightarrow C\) be an epimorphism,
where \(C\) is connected and all full fibres \(F_a=\pi^{-1}(a)\)
belong to \(\OS\). Then \(E\in\OS\).
\end{lemma}
\begin{proof}
Choose an arbitrary first orbit \(O=\Orb(x,E)\).
By Lemma~\ref{lem:images}, \(\pi(O)=C\), and by
Lemma~\ref{lem:slices}, \(O\cap F_a\in\OS\) for every \(a\in C\).

Set \(G=\Inn(E)\), \(H=\Inn(C)\),
\(c=L_x^E\), \(h=L_{\pi(x)}^C\).
The epimorphism \(\pi\) induces \(\sigma:G\twoheadrightarrow H\).
The map \(\lambda_E(y)=L_y^E\) is a homomorphism of
quandles, since
\[
 L_{y\conj z}^E=L_y^E L_z^E(L_y^E)^{-1}.
\]
Equivariance of this map gives
\(\lambda_E(O)=c^G\).
Connectedness of \(C\) similarly gives
\[
 \{L_a^C:a\in C\}=h^H,\qquad H=\langle h^H\rangle.
\]

Write \(b=\pi(x)\) and
\[
 T=\{a\in C:L_a^C=L_b^C\},\qquad
 S=O\cap\pi^{-1}(T).
\]
The quandle \(T\) is trivial: for \(a,d\in T\),
\(a\conj d=L_a^C(d)=L_d^C(d)=d\).
The map \(S\twoheadrightarrow T\) has full fibres
\(O\cap F_a\in\OS\); by Lemma~\ref{lem:trivialbase},
\(S\in\OS\).

The full fibre of \(c^G\to h^H\) over \(h\) is exactly
\(\lambda_E(S)\). Indeed, the image of an element
\(L_y^E\), with \(y\in O\), is \(h\) if and only if
\(L_{\pi(y)}^C=L_b^C\).
By Lemma~\ref{lem:images}, this fibre belongs to \(\OS\).
Proposition~\ref{prop:groupquandle} gives \(c^G\in\OS\).

The epimorphism \(\lambda_E|_O:O\twoheadrightarrow c^G\)
is a covering: equality of global translations
\(L_y^E=L_z^E\) implies equality of their restrictions
\(L_y^O=L_z^O\).
By Lemma~\ref{lem:covering}, \(O\in\OS\).
This holds for every first orbit in \(E\);
therefore, every orbit series in \(E\) stabilizes.
\end{proof}

\begin{proof}[Proof of Theorem~\ref{thm:main}]
Take an arbitrary series \(Q_i\) in \(Q\).
The images \(B_i=\pi(Q_i)\) form an orbit series in \(B\)
and from some index \(n\) onward are all equal to a connected subquandle \(B_*\).
The restriction \(Q_n\twoheadrightarrow B_*\) has full fibres
\(Q_n\cap F_b\in\OS\) by Lemma~\ref{lem:slices}.
Applying Lemma~\ref{lem:connectedbase}, we obtain \(Q_n\in\OS\).
Therefore the chosen tail of the series stabilizes.
Since the original series was arbitrary, \(Q\in\OS\).
\end{proof}

\begin{remark}
A quandle is called \emph{faithful} if \(y\mapsto L_y\)
is injective \cite[Section~2.2]{BCNW}.
The map \(y\mapsto L_y\) need not be injective.
Thus the proof also covers nonfaithful quandles.
No common finite bound on the depths of the full fibres is needed.
The bound in Theorem~\ref{thm:group} concerns group series;
an optimal quantitative bound for arbitrary
quandles is not established in this work.
\end{remark}

\section{\texorpdfstring{Negative answer for \(\OS_\omega\)}
{Negative answer for OS omega}}\label{sec:negative}

\begin{proof}[Proof of Proposition~\ref{prop:negative}]
Let \(D_{2^r}\) be the dihedral quandle on \(\Z/2^r\Z\)
with operation \(a\conj b=2a-b\).
Consider the disjoint union
\[
 Q=\bigsqcup_{r\ge1}D_{2^r},
\]
where, for elements in different summands, \(x\conj y=y\).
This is the quandle in \cite[Example~4.2]{BCNW}.
The projection \(\pi:Q\to T_{\mathbb N_{\ge1}}\)
onto the summand index is an epimorphism to a trivial quandle.
Its full fibres are finite \(D_{2^r}\), so they all
belong to \(\OS_\omega\); the base belongs to \(\OS_1\).

The depths in the domain are unbounded.
Choose zero in the summand \(D_{2^r}\) and use the constant anchor \(0\).
The inner orbits of an even dihedral quandle
are the parity classes: \(L_aL_b\) is translation by \(2(a-b)\).
Elements of the other summands act trivially.
Thus the series in \(Q\) has the form
\[
 Q_0=Q,\qquad
 Q_i=2^i\Z/2^r\Z\quad(1\le i\le r),\qquad
 Q_r=Q_{r+1}=\{0\}.
\]
Every step up to \(r\) is strict.
For any proposed common bound \(n\), choose \(r>n\);
then \(Q\notin\OS_n\).
Consequently, \(Q\notin\OS_\omega\).

For any coset \(S=a+2^j\Z/2^r\Z\) and \(x\in S\),
we have
\[
 \Orb(x,S)=x+2^{j+1}\Z/2^r\Z.
\]
Indeed, the values \(L_y(x)=2y-x\) for \(y\in S\)
run through the entire indicated coset, which is invariant
under all \(L_y^{\pm1}\).
Therefore every individual series, after its first step, lies
in one summand and, for any subsequent choice of anchors,
reaches a singleton quandle by step \(r\) at the latest.
Hence \(Q\in\tOS\subseteq\OS\), in agreement with
Theorem~\ref{thm:main}.
\end{proof}

\begin{remark}\label{rmk:uniform}
Each fibre in the last example has finite depth, but there is no common
index for all fibres. If one defines a different notion
of closure requiring a common bound on depths, this example does not refute it.
The same example also shows that the class
\(\tOS_\omega=\bigcup_{n\ge1}\tOS_n\) is not closed under extensions:
the fibres belong to \(\tOS_\omega\), while the domain does not.
Consequently, the assertion that \(\tOS_\omega\) is closed under extensions
in \cite[Corollary~5.6, arXiv v3]{BCNW}
is false for the definition of extension in its Section~2.2,
which does not require a uniform bound across fibres.
\cite[Lemma~5.5, arXiv v3]{BCNW} assumes a single fixed
index \(m\) for all full fibres and is not contradicted by this example.
This hypothesis must be retained to deduce membership in \(\tOS_\omega\)
from that lemma.
\end{remark}

\begin{thebibliography}{9}
\raggedright
\bibitem{BCNW}
M.~Bonatto, A.~S.~Crans, T.~Nasybullov, G.~T.~Whitney.
\newblock Quandles with orbit series conditions.
\newblock \emph{Journal of Algebra}, \textbf{567} (2021), 284--309.
\newblock DOI:
\href{https://doi.org/10.1016/j.jalgebra.2020.09.026}
{10.1016/j.jalgebra.2020.09.026}.
\newblock The numbering in the text follows
\href{https://arxiv.org/abs/2004.10227v3}{arXiv:2004.10227v3},
August 18, 2020.
\bibitem{Eisermann}
M.~Eisermann.
\newblock Quandle coverings and their Galois correspondence.
\newblock \emph{Fundamenta Mathematicae}, \textbf{225} (2014), 103--167.
\newblock DOI:
\href{https://doi.org/10.4064/fm225-1-7}{10.4064/fm225-1-7}.
\newblock Numbering:
\href{https://arxiv.org/abs/math/0612459v3}{arXiv:math/0612459v3}.
\end{thebibliography}

\end{document}
